\subsection{Implication}

\newcommand{\norimplies}[2]{\norvee{\norneg{#1}}{#2}} % NOR hacking
\begin{notation}
    \[P \implies Q \coloneqq \neg P \vee Q \qquad [\norimplies{P}{Q}]\]
\end{notation}

% [P]...Q > P=>Q
\begin{align*}
    \premise{[P] \cdots Q}\\
    \assumption[1]{\neg P \nor Q}\\
    \assumption[2]{P}\\
    \byrule[2]{Q}{\text{given}}\\
    \byrule[2]{\bot}{\ref{rule:nor-elimination}}\\
    \byrule[1]{\neg P}{\ref{rule:negation-introduction}}\\
    \byrule[1]{\bot}{\ref{rule:nor-elimination}}\\
    \byrule{\neg(\neg P \nor Q)}{\ref{rule:negation-introduction}}\\
    \bydefinition{\neg P \vee Q}\\
    \bydefinition{P \implies Q}
\end{align*}
\begin{equation}
    \inferrule{[P] \cdots Q}{P \implies Q}
    \tag{{$\implies$}I}
    \label{rule:conditional-introduction}
\end{equation}

% P, P=>Q > Q
\begin{align*}
    \premise{P}\\
    \premise{P \implies Q}\\
    \bydefinition{\neg(\neg P \nor Q)}\\
    \assumption[1]{\neg Q}\\
    \assumption[2]{\neg P}\\
    \byrule[2]{\bot}{\ref{rule:negation-elimination}}\\
    \assumption[2]{Q}\\
    \byrule[2]{\bot}{\ref{rule:negation-elimination}}\\
    \byrule[1]{\neg P \nor Q}{\ref{rule:nor-introduction}}\\
    \byrule[1]{\bot}{\ref{rule:negation-elimination}}\\
    \byrule{Q}{\ref{rule:reductio-ad-absurdum}}
\end{align*}
\begin{equation}
    \inferrule{P \\ P \implies Q}{Q}
    \tag{{$\implies$}E}
    \label{rule:conditional-elimination}
\end{equation}

% > P => P
\begin{align*}
    \assumption[1]{P}\\
    \given[1]{P}\\
    \byrule{P \implies P}{\ref{rule:conditional-introduction}}    
\end{align*}
\begin{equation}
    \vdash P \implies P
    \label{rule:conditional-reflexivity}
\end{equation}
